\section{讨论}
\label{discussion:section}

相位表示将全局正交性与局部几何分开处理。酉性保证完整测量，而式~\eqref{core:signal} 将每个带符号局部分量分解为背景期望和涉及物理零串修正的干涉。四特征核对全部 $n\ge5$ 消去背景，使这一干涉控制局部极化。随维数变化的字符和归约为同一组有限消去条件，因此该机制同时适用于奇数和偶数比特。$t=\pi/2$ 时单点修正消失，$t=-\pi/2$ 时其干涉使基准信号加倍。通过式~\eqref{core:spectrum} 中共同的约化谱，这种相位控制对每个比特、每个测量结果产生相同的单比特对其余系统纠缠变化：半径增加时，该纠缠下降。

可容许路径的形状同样影响增强能力。四律分类约束两值相位直线，而 $n=4$ 的结果表明，即使允许任意实直线生成元，仍可能存在可由弯曲路径突破的障碍。在这一情形，高阶正四面体条件将可容许直线方向限制在两个平面内，增强路径则离开这两个平面。两个相位方向同时存在，也不足以保证它们能够独立变化。在偶数比特构造中，这一兼容性取决于混合背景项。对偶数 $n\ge8$，至少两个奇数内部位允许保持奇偶性的双位翻转，剩余的固定规模和式则通过核的性质相消。$n=6$ 时，唯一的奇数内部位留下正比于 $\sin\alpha\sin\beta$ 的项，阻碍两个方向的直接叠加。完整四颜色相位模型在基准点附近也只保留两条相交分支。

这一几何结构引出两个问题。在 $n=4$ 情形，文献~\cite{PauwelsGisin2026} 中大半径族的指定参数段不包含与本文基准 EJM 在局域酉变换、比特置换和结果重标号下等价的测量。本文构造的弯曲分支能否在正四面体相位空间的其他位置与该族连通，仍未解决。在 $n=6$ 情形，更细的相位划分可能允许同时包含两支的双参数族；判断其是否存在，需要控制四颜色相位模型之外新增的混合项。
